%! Linear Functions: Slope, Forms & Graphing — Unit 1, Lesson 1.1 — Experience & Formalize
%! (Activity · QuickNotes · Application · Check Your Understanding) (Answer Key)
% Directory stays `experience/` (build identifier in shared/lesson.mk); the label the student and
% the teacher read is "Experience & Formalize".
% Math Medic sizing: 12pt body, OPEN white space after every question (no writing lines).
% Page budget (user, 2026-08-20): Activity 2pp max · QuickNotes 1/2pp · Application 1/2-1pp ·
% Check Your Understanding 1-2pp (practice, NOT scored — there is no homework).
\documentclass[12pt]{article}
\usepackage{algebra2-article}
\usepackage{algebra2-key}
\newcommand{\answerspace}[2]{\par\nopagebreak\noindent\begin{minipage}[t][#1][t]{\linewidth}%
  \color{keyred}\bfseries #2\end{minipage}\par}
\newcommand{\lgrid}[4]{%
  \draw[step=1, linegray!40, very thin] (#1,#3) grid (#2,#4);
  \draw[->, semithick, charcoal] (#1,0)--(#2,0) node[below right, font=\scriptsize]{$x$};
  \draw[->, semithick, charcoal] (0,#3)--(0,#4) node[left, font=\scriptsize]{$y$};}
% Gym A's cost graph: climbs n on 0..8, cost C on 0..30 (ticks every $6).
\newcommand{\gymgrid}{%
  \draw[xstep=1, ystep=3, linegray!40, very thin] (0,0) grid (8,30);
  \draw[->, semithick, charcoal] (0,0)--(8.7,0) node[right, font=\scriptsize]{$n$};
  \draw[->, semithick, charcoal] (0,0)--(0,32) node[above, font=\scriptsize]{$C$ (\$)};
  \foreach \x in {2,4,6,8} \node[charcoal, font=\tiny, below] at (\x,0){\x};
  \foreach \y in {6,12,18,24,30} \node[charcoal, font=\tiny, left] at (0,\y){\y};}

\begin{document}

\pageheader{Unit 1, Lesson 1.1}{Experience \& Formalize: Two Receipts --- Answer Key}

\vspace{0.06in}

\begin{headlinebox}{forestbg}
\textbf{Two Receipts.} Two climbing gyms charge the same way --- something to walk in the door, then a
price per climb --- but they tell you about it very differently. Work both with your group using only
what you already know, and be ready to show \emph{where} each number came from.
\end{headlinebox}

\vspace{0.04in}

% ══ Part 1 of 4 — ACTIVITY (20 min, groups of 3–4, prior knowledge only) ══════
\begin{scenariobox}[1. Summit Gym --- They Show You the Graph]{forest}
\begin{minipage}[t]{0.54\linewidth}
\vspace{0pt}
\begin{enumerate}[label=\textbf{\alph*.}, leftmargin=1.8em, itemsep=4pt]
  \item What does a visit cost if you climb \emph{zero} times? \$\ans{6} \;
        What is that charge, in the story?
        \answerspace{1.2cm}{The walk-in (entry) fee --- what you pay before climbing anything.}
  \item Each extra climb adds \$\ans{3}. Mark on the graph where you see that, and say how
        you found it.
        \answerspace{1.4cm}{From $(0,6)$ to $(4,18)$ the cost rises \$12 over 4 climbs, so \$3 per climb --- and every single step right goes up 3.}
  \item Write one rule for the cost $C$ of $n$ climbs. \; $C(n) = $ \ans{$3n + 6$}
\end{enumerate}
\end{minipage}\hfill
\begin{minipage}[t]{0.42\linewidth}
\vspace{0pt}\centering
\begin{tikzpicture}[x=0.52cm, y=0.13cm]
  \gymgrid
  \draw[very thick, forest] (0,6)--(8,30);
  \fill[charcoal] (0,6) circle (2.6pt);
  \fill[charcoal] (4,18) circle (2.6pt) node[right, font=\tiny, charcoal]{$(4,18)$};
\end{tikzpicture}\\[1pt]
{\footnotesize Summit Gym: cost of $n$ climbs}
\end{minipage}

\begin{enumerate}[label=\textbf{\alph*.}, start=4, leftmargin=1.8em, itemsep=4pt]
  \item Use your rule to price 7 climbs: \$\ans{27}. Check it against the graph --- do the
        rule and the picture agree? \ans{yes}
\end{enumerate}
\end{scenariobox}

\boxguard[14]
\begin{scenariobox}[2. Basecamp Gym --- All You Get Is Two Receipts]{forest}
No graph, no posted prices. A friend hands you two receipts: \textbf{3 climbs cost \$24}, and
\textbf{7 climbs cost \$32}.

\begin{enumerate}[label=\textbf{\alph*.}, leftmargin=1.8em, itemsep=4pt]
  \item Going from the first receipt to the second is \ans{4} more climbs for
        \$\ans{8} more, so each climb adds \$\ans{2}.
  \item Work back to what Basecamp charges for \emph{zero} climbs.
    \begin{work}
      \text{3 climbs} &= \$24 \\
      \text{3 climbs alone cost } 3(\$2) &= \$6 \\
      \text{so the walk-in fee is } \$24 - \$6 &= \$18
    \end{work}
    Walk-in fee: \$\ans{18}
  \item Write one rule for Basecamp. \; $D(n) = $ \ans{$2n + 18$}
  \item \textbf{Priya never found the walk-in fee.} She started from the receipt she had and wrote
        $$D - 24 = 2(n - 3).$$
        Is her rule the same as yours, or a different one? Test both at $n = 7$, then tidy hers up.
    \begin{work}
      D - 24 &= 2(n-3) \\
      D - 24 &= 2n - 6 \\
           D &= 2n + 18
    \end{work}
    \answerspace{1.6cm}{The same rule. At $n=7$ both give \$32, and expanding Priya's gives $D = 2n+18$ exactly. Hers is built from a receipt she had; ours is built from the walk-in fee.}
  \item Cheaper for \textbf{2 climbs}? \ans{Summit} \; For \textbf{15}?
        \ans{Basecamp} \; What makes the answer switch somewhere?
    \answerspace{1.4cm}{Basecamp starts \$12 higher but grows \$1 slower per climb, so the gap closes as $n$ grows and the cheaper gym changes.}
  \item If you drew Basecamp on the same axes as Summit, it would cross the vertical axis at
        \ans{$18$} and be \ans{less steep} than Summit. (Two numbers --- no drawing.)
\end{enumerate}
\end{scenariobox}

% ══ Part 2 of 4 — QUICKNOTES (13 min debrief fills this; target half a page) ══
\boxguard[14]
\begin{tcolorbox}[enhanced, colback=sky, colframe=navy, boxrule=0.8pt, arc=2mm,
    left=3mm, right=3mm, top=1.5mm, bottom=1.5mm, breakable,
    fonttitle=\bfseries\small\color{white}, title={QuickNotes: Three Ways to Write One Line},
    attach boxed title to top left={yshift=-2mm, xshift=4mm},
    boxed title style={colback=navy, colframe=navy, arc=1.5mm}]
\small
\begin{minipage}[t]{0.30\linewidth}
\vspace{2pt}\centering
\begin{tikzpicture}[scale=0.26]
  \lgrid{-4}{5}{-6}{6}
  \draw[very thick, navy, dashed] (-4,-4)--(5,5) node[above left, font=\tiny, navy]{$y=x$};
  \draw[very thick, forest] (-1,-6)--(5,6) node[below right, font=\tiny, forest]{$y=2x-4$};
  \fill[charcoal] (0,-4) circle (3.6pt);
  \fill[charcoal] (2,0) circle (3.6pt);
\end{tikzpicture}\\[1pt]
{\footnotesize The parent $y=x$, stretched and shifted}
\end{minipage}\hfill
\begin{minipage}[t]{0.66\linewidth}
\vspace{0pt}
\begin{itemize}[leftmargin=1.1em, itemsep=2pt, topsep=0pt]
  \item \textbf{Slope} from any two points: $m = \dfrac{y_2-y_1}{x_2-x_1} = \dfrac{\Delta y}{\Delta x}$.
        Any two work because the rate is \ans{constant}.
  \item \textbf{Slope-intercept} $y = mx+b$ --- hands you the \ans{rate} and the
        \ans{starting value}. (Your Summit rule.)
  \item \textbf{Point-slope} $y - y_1 = m(x - x_1)$ --- built from a slope and
        \ans{any point at all}. (Priya's rule.)
  \item \textbf{Standard} $Ax + By = C$, integers, $A \ge 0$ --- tidy for reading
        \ans{intercepts}.
\end{itemize}
\end{minipage}

\vspace{3pt}
\begin{itemize}[leftmargin=1.1em, itemsep=2pt, topsep=0pt]
  \item \textbf{Same line, three outfits.} Expand point-slope $\to$ slope-intercept; clear fractions
        and move $x$ across $\to$ standard. \emph{Different-looking equations are not different
        lines} --- that was Priya.
  \item \textbf{Transforming the parent} $y = x$: multiplying by $k$ changes the
        \ans{slope} (steeper, flatter, or flipped); adding $k$ changes the
        \ans{$y$-intercept}. Every non-vertical line is $y=x$ after those two moves.
  \item \textbf{Parallel} lines have \ans{equal} slopes; \textbf{perpendicular} slopes are
        \ans{opposite reciprocals} ($m_1m_2 = -1$). \; Horizontal $y=c$: slope
        \ans{$0$}. \; Vertical $x=a$: slope \ans{undefined}.
\end{itemize}
\end{tcolorbox}

% ══ Part 3 of 4 — APPLICATION (7 min, worked together; the modeling standard) ══
\boxguard[14]
\begin{notesbox}{Application: The Drone Descent}
We will do this one together. A drone comes down at a steady rate. At $t = 2$ seconds it is
$46$~m up; at $t = 6$ seconds it is $30$~m up. Let $h$ be its height after $t$ seconds.

\begin{enumerate}[leftmargin=1.7em, itemsep=4pt]
  \item Find the rate.
    \begin{work}
      m &= \frac{30 - 46}{6 - 2} \\
        &= \frac{-16}{4} \\
        &= -4
    \end{work}
    So the drone falls \ans{4 m each second}.
  \item Write the line in \textbf{point-slope} form using $(2,46)$: \ans{$h-46=-4(t-2)$}
        \\[2pt] Now put it in \textbf{slope-intercept} form.
    \begin{work}
      h - 46 &= -4(t - 2) \\
      h - 46 &= -4t + 8 \\
           h &= -4t + 54
    \end{work}
    The $54$ means \ans{it started at 54 m} \; and in \textbf{standard} form the line is
    \ans{$4t+h=54$}
  \item Suppose the drone had begun the same descent from $60$~m instead. In $h = mt + b$, which
        number changes and which stays put --- and what does that do to the graph?
    \answerspace{1.2cm}{$b$ becomes $60$; $m$ stays $-4$. The line starts higher but falls at the same steepness --- a shift up, not a stretch.}
\end{enumerate}
\end{notesbox}

% ══ Part 4 of 4 — CHECK YOUR UNDERSTANDING (10 min, practice, NOT scored) ═════
% The lesson's only practice set — there is no homework — so the six items carry the whole
% A.F.1a–e spread: slope, the three forms and converting, transformations of the parent,
% the special cases, a model in a fresh context, and the formative MC on perpendicular slopes.
\boxguard[7]
\begin{notesbox}{Check Your Understanding \quad {\normalfont\itshape (practice --- not scored)}}
Six exercises --- the lesson's practice, and there is \textbf{no homework}. Work 1--3 with a partner,
then 4--6 on your own.

\begin{enumerate}[leftmargin=1.7em, itemsep=3pt]
  \item \textbf{Slope from two points.} \;
        Through $(-2,7)$ and $(4,-5)$: $m = {\color{keyred}\mathbf{-2}}$ \quad
        Through $(-3,4)$ and $(6,4)$: $m = {\color{keyred}\mathbf{0}}$ \quad
        Through $(5,1)$ and $(5,-3)$: $m = $ \ans{undefined}

  \boxguard[12]
  \item \textbf{One line, two ways.} The line shown passes through $(0,-1)$ and $(3,5)$.
    \begin{center}
    \begin{minipage}[c]{0.32\linewidth}\centering
    \begin{tikzpicture}[scale=0.26]
      \lgrid{-3}{5}{-6}{7}
      \draw[very thick, forest] (-2.5,-6)--(4,7);
      \fill[charcoal] (0,-1) circle (4pt);
      \fill[charcoal] (3,5) circle (4pt);
    \end{tikzpicture}
    \end{minipage}\hfill
    \begin{minipage}[c]{0.64\linewidth}
    \begin{enumerate}[label=(\alph*), leftmargin=1.9em, itemsep=3pt]
      \item Slope: \ans{$2$} \; Slope-intercept form: \ans{$y=2x-1$}
      \item Point-slope form, using $(3,5)$: \ans{$y-5=2(x-3)$}
      \item Are (a) and (b) the same line? \ans{yes} \; Show it:
    \end{enumerate}
    \end{minipage}
    \end{center}
    \begin{work}
      y - 5 &= 2(x - 3) \\
      y - 5 &= 2x - 6 \\
          y &= 2x - 1
    \end{work}

  \item \textbf{Change its outfit.} Rewrite $y + 2 = -\tfrac34(x - 4)$ in slope-intercept form,
        then in standard form.
    \begin{work}
      y + 2 &= -\tfrac34(x-4) \\
      y + 2 &= -\tfrac34 x + 3 \\
          y &= -\tfrac34 x + 1 \\
         4y &= -3x + 4 \\
    3x + 4y &= 4
    \end{work}
    Slope-intercept: \ans{$y=-\tfrac34x+1$} \quad Standard: \ans{$3x+4y=4$}

  \boxguard[12]
  \item \textbf{The parent, and the two odd ones.} Start from $y = x$ and end at $y = -3x + 2$.
    \begin{enumerate}[label=(\alph*), leftmargin=1.9em, itemsep=3pt]
      \item Steeper or flatter? \ans{steeper} \; Flipped? \ans{yes} \;
            Shifted? \ans{up 2}
      \item Slope of $y = 4$: \ans{$0$} \quad Slope of $x = 4$: \ans{undefined}
      \item Which of those two \emph{cannot} be written as $y = mx + b$, and why?
        \answerspace{1.2cm}{$x=4$. There is no slope to put in for $m$ --- the run between any two of its points is $0$, and one $x$ gives every $y$, so it is not a function of $x$ at all.}
    \end{enumerate}

  \boxguard[12]
  \item \textbf{Model it.} A candle burns at a steady rate. After $1$ hour it is $17$~cm tall;
        after $4$ hours it is $11$~cm tall.
    \begin{enumerate}[label=(\alph*), leftmargin=1.9em, itemsep=3pt]
      \item Rate: \ans{$-2$ cm per hour} \; Rule: $h(t) = $ \ans{$-2t + 19$}
      \item The $y$-intercept means \ans{it began at 19 cm} \; After $6$ hours:
            \ans{7 cm} \; Burns out at \ans{$t=9.5$ h}
    \end{enumerate}

  \item \textbf{Multiple choice.} Which line is \emph{perpendicular} to $y = \tfrac12 x - 4$ and
        passes through $(2,-1)$?
    \begin{enumerate}[label=\Alph*., leftmargin=2.2em, itemsep=1pt]
      \item[\textbf{A.}] $y = \tfrac12 x - 2$ \hfill \textbf{B.} $y = -2x + 3$ \textcolor{keyred}{\textbf{$\leftarrow$ correct}}
      \item[\textbf{C.}] $y = 2x - 5$ \hfill \textbf{D.} $y = -\tfrac12 x$
    \end{enumerate}
    Say in one sentence how you ruled out one of the others.
    \answerspace{1.2cm}{B: perpendicular needs the opposite reciprocal of $\tfrac12$, which is $-2$, and $-1 = -2(2)+3$ checks. (A is parallel, not perpendicular; C flips the sign but not the reciprocal; D takes the reciprocal but not the sign.)}
\end{enumerate}
\end{notesbox}

\boxguard[5]
\begin{spiralbox}
\textbf{Coming next --- Lesson 1.2: Absolute Value Equations \& Inequalities.} Every line you wrote
today crosses zero once. Next comes a rule that reaches the same value \emph{twice}, because it
measures \emph{distance} --- so one equation will have two answers.
\end{spiralbox}

\end{document}
